%----------------------------------------------------------------------------------------
% NOTATIE EN FIGUREN VAN DE CURSUS KEGELSNEDEN
% Staat in de preambule en niet in een hoofdstuk: ook de lijsten met oplossingen gebruiken dit.
%----------------------------------------------------------------------------------------
% Bewaarde grafiek, geplaatst door maak-pdf.py: bestand en pagina in figuren-cache/
\newcommand{\grafiekbestand}[2]{\includegraphics[page=#2,trim=3bp 3bp 3bp 3bp]{figuren-cache/#1}}

\providecommand{\R}{\mathbb{R}}
\newcommand{\dx}{\,\mathrm{d}x}
\newcommand{\e}{\mathrm{e}}
% Afstand van een punt tot een punt of een rechte: \afst{P}{l} geeft d(P, l)
\newcommand{\afst}[2]{d(#1, #2)}

% Zoals stappen, maar met een bredere kolom voor de berekening: voor lange formules
\newenvironment{stappenbreed}{%
  \par\renewcommand{\arraystretch}{1.25}%
  \noindent\begin{tabular}{@{}>{$\displaystyle}p{9.8cm}<{$}@{\hspace{4mm}}>{\RaggedRight\small\itshape\arraybackslash}p{4.9cm}@{}}%
}{\end{tabular}\par}

% Zoals stappen, met eigen kolombreedtes: \begin{stappenvrij}{6.2cm}{3.8cm}. Voor naast een figuur.
\newenvironment{stappenvrij}[2]{%
  \par\renewcommand{\arraystretch}{1.25}%
  \noindent\begin{tabular}{@{}>{$\displaystyle}p{#1}<{$}@{\hspace{3mm}}>{\RaggedRight\small\itshape\arraybackslash}p{#2}@{}}%
}{\end{tabular}\par}

% Tekst of berekening links, figuur rechts: \naast[breedte rechts]{links}{rechts}
% De breedte rechts komt uit de figuur zelf; het optionele argument wordt niet gebruikt.
\newsavebox{\naastbox}
\newcommand{\naast}[3][]{%
  \par\noindent\sbox{\naastbox}{#3}%
  \begin{minipage}[t]{\dimexpr\linewidth-\wd\naastbox-3mm\relax}\vspace{0pt}#2\end{minipage}\hfill
  \begin{minipage}[t]{\wd\naastbox}\vspace{0pt}\usebox{\naastbox}\end{minipage}\par\addvspace{1mm}}
% Een rij figuren, gecentreerd
\newcommand{\figuren}[1]{\par\addvspace{1mm}{\centering #1\par}\addvspace{1.5mm}}
% Opgave met een grafiek rechts: \begin{metgrafiek}[breedte]{tikz} opgave en deelvragen \end{metgrafiek}
% De solution komt na \end{metgrafiek}, zoals bij metprent.
\newsavebox{\grafiekbox}
\newenvironment{metgrafiek}[2][]{%
  \sbox{\grafiekbox}{#2}%
  \begin{minipage}[t]{\dimexpr\linewidth-\wd\grafiekbox-3mm\relax}%
}{%
  \end{minipage}\hfill
  \begin{minipage}[t]{\wd\grafiekbox}\vspace{-2mm}\usebox{\grafiekbox}\end{minipage}\par}

% Rooster met assen: \rooster{xmin}{xstap}{xmax}{ymin}{ystap}{ymax}, met getallen op de assen (gehele waarden).
% \roosterleeg doet hetzelfde zonder getallen, voor eigen labels.
\newcommand{\roosterleeg}[6]{%
  \draw[black!14,very thin,xstep=#2,ystep=#5] (#1,#4) grid (#3,#6);
  \draw[->,grijs] (#1,0) -- ({#3+.4*#2},0) node[right,text=black] {$x$};
  \draw[->,grijs] (0,#4) -- (0,{#6+.4*#5}) node[above,text=black] {$y$};}
\newcommand{\rooster}[6]{%
  \roosterleeg{#1}{#2}{#3}{#4}{#5}{#6}
  \pgfmathtruncatemacro{\rtweex}{#1+#2}\pgfmathtruncatemacro{\rtweey}{#4+#5}
  \foreach \t in {#1,\rtweex,...,#3}{\ifnum\t=0\else\node[below,inner sep=1.5pt,text=black] at (\t,0) {$\t$};\fi}
  \foreach \t in {#4,\rtweey,...,#6}{\ifnum\t=0\else\node[left,inner sep=1.5pt,text=black] at (0,\t) {$\t$};\fi}}

% In de functievoorschriften is \x de veranderlijke; schrijf (\x)^2 en sin(\x r).
\newcommand{\assen}[4]{%
  \draw[->,lijngrijs] (#1,0) -- (#2,0) node[right,text=black] {$x$};
  \draw[->,lijngrijs] (0,#3) -- (0,#4) node[above,text=black] {$y$};}
\newcommand{\kromme}[4][]{\draw[accent,thick,#1] plot[domain=#2:#3,samples=90,smooth] (\x,{#4});}
% \gebiedtussen: tussen twee grafieken (voor oppervlaktes, zoals in Integralen)
\newcommand{\gebiedtussen}[5][accent!30]{\fill[#1] plot[domain=#2:#3,samples=70] (\x,{#4}) -- plot[domain=#3:#2,samples=70] (\x,{#5}) -- cycle;}

% KEGELSNEDEN. Optioneel argument: extra opties voor de lijn (dashed, very thick, een kleur ...).
% \ellips{x0}{y0}{a}{b}: middelpunt (x0, y0), halve assen a (horizontaal) en b (verticaal)
\newcommand{\ellips}[5][]{\draw[accent,thick,#1] (#2,#3) ellipse ({#4} and {#5});}
% \hyperbool{x0}{y0}{a}{b}{tmax}: brandpunten op een horizontale lijn, beide takken, parameter t van -tmax tot tmax
\newcommand{\hyperbool}[6][]{%
  \draw[accent,thick,#1] plot[domain=-#6:#6,samples=60,smooth] ({#2+(#4)*cosh(\x)},{#3+(#5)*sinh(\x)});
  \draw[accent,thick,#1] plot[domain=-#6:#6,samples=60,smooth] ({#2-(#4)*cosh(\x)},{#3+(#5)*sinh(\x)});}
% \hyperboolv{x0}{y0}{a}{b}{tmax}: brandpunten op een verticale lijn; a is de halve afstand tussen de toppen
\newcommand{\hyperboolv}[6][]{%
  \draw[accent,thick,#1] plot[domain=-#6:#6,samples=60,smooth] ({#2+(#5)*sinh(\x)},{#3+(#4)*cosh(\x)});
  \draw[accent,thick,#1] plot[domain=-#6:#6,samples=60,smooth] ({#2+(#5)*sinh(\x)},{#3-(#4)*cosh(\x)});}
% \parabool{x0}{y0}{p}{ymin}{ymax}: (y - y0)^2 = 2p(x - x0), opening naar rechts als p > 0; y loopt van ymin tot ymax
\newcommand{\parabool}[6][]{\draw[accent,thick,#1] plot[domain=#5:#6,samples=60,smooth] ({#2+((\x)-(#3))^2/(2*(#4))},{\x});}
% \paraboolv{x0}{y0}{p}{xmin}{xmax}: (x - x0)^2 = 2p(y - y0), opening naar boven als p > 0
\newcommand{\paraboolv}[6][]{\draw[accent,thick,#1] plot[domain=#5:#6,samples=60,smooth] ({\x},{#3+((\x)-(#2))^2/(2*(#4))});}
% Een punt als stip, en een brandpunt (iets groter) met een label: \stip{x}{y}, \brandpunt[positie]{x}{y}{label}
\newcommand{\stip}[2]{\fill (#1,#2) circle (1.5pt);}
\newcommand{\brandpunt}[4][below]{\fill[accent] (#2,#3) circle (1.9pt) node[#1,text=black,inner sep=2pt] {$#4$};}
% Een rechte als lijn door twee punten, verlengd: \rechte[opties]{x1}{y1}{x2}{y2}
\newcommand{\rechte}[5][]{\draw[lijngrijs,thick,shorten <=-2mm,shorten >=-2mm,#1] (#2,#3) -- (#4,#5);}

% DUBBELKEGEL in een ruimtelijke figuur: optie [kegel] bij tikzpicture. De as is de z-as, de top de oorsprong,
% tan(halve tophoek) = 0,6 en de kegel loopt van z = -2 tot z = 2. De y-as wijst naar de kijker (naar onder op papier).
\tikzset{kegel/.style={x={(.72cm,0cm)},y={(0cm,-.2cm)},z={(0cm,.72cm)}}}
\newcommand{\dubbelkegel}{%
  \draw[lijngrijs] plot[domain=0:360,samples=60,smooth,variable=\t] ({1.2*cos(\t)},{1.2*sin(\t)},2);
  \draw[lijngrijs] plot[domain=0:360,samples=60,smooth,variable=\t] ({1.2*cos(\t)},{1.2*sin(\t)},-2);
  \draw[lijngrijs] (-1.2,0,2) -- (1.2,0,-2) (1.2,0,2) -- (-1.2,0,-2);
  \draw[lijngrijs,densely dotted] (0,0,-2.5) -- (0,0,2.5);
  \fill (0,0,0) circle (1.2pt);}
% Een snijvlak z = m x + h als parallellogram: \snijvlak{m}{h}{xmin}{xmax}
\newcommand{\snijvlak}[4]{\fill[accent!25,opacity=.55] (#3,-1.5,{#1*(#3)+#2}) -- (#4,-1.5,{#1*(#4)+#2}) -- (#4,1.5,{#1*(#4)+#2}) -- (#3,1.5,{#1*(#3)+#2}) -- cycle;
  \draw[accent!60] (#3,-1.5,{#1*(#3)+#2}) -- (#4,-1.5,{#1*(#4)+#2}) -- (#4,1.5,{#1*(#4)+#2}) -- (#3,1.5,{#1*(#3)+#2}) -- cycle;}
% De doorsnede met dat vlak: \kegelsnede{m}{h}{hoek van}{hoek tot}, voor een vlak dat de as boven de top snijdt
\newcommand{\kegelsnede}[4]{\draw[accent,very thick] plot[domain=#3:#4,samples=80,smooth,variable=\t]
  ({.6*#2/(1-.6*(#1)*cos(\t))*cos(\t)},{.6*#2/(1-.6*(#1)*cos(\t))*sin(\t)},{#2/(1-.6*(#1)*cos(\t))});}
% Een verticaal snijvlak y = c en de hyperbool die het uitsnijdt: \hyperboolsnede{c}{umax}
\newcommand{\hyperboolsnede}[2]{%
  \fill[accent!25,opacity=.55] (-1.8,#1,-2.3) -- (1.8,#1,-2.3) -- (1.8,#1,2.3) -- (-1.8,#1,2.3) -- cycle;
  \draw[accent!60] (-1.8,#1,-2.3) -- (1.8,#1,-2.3) -- (1.8,#1,2.3) -- (-1.8,#1,2.3) -- cycle;
  \draw[accent,very thick] plot[domain=-#2:#2,samples=60,smooth,variable=\t] ({#1*sinh(\t)},#1,{#1/.6*cosh(\t)});
  \draw[accent,very thick] plot[domain=-#2:#2,samples=60,smooth,variable=\t] ({#1*sinh(\t)},#1,{-#1/.6*cosh(\t)});}
